%--------------------------------------------------------------------------------
%
%
%--------------------------------------------------------------------------------
\chapter*{Addressing and Access Control}
\addcontentsline{toc}{chapter}{Addressing and Access Control}

The CPU emits a 52-bit virtual addresses. Each virtual memory access needs to be translated and checked for valid access rights.

\section*{Address Translation}
\addcontentsline{toc}{section}{Address Translation}

Each virtual memory access generated by the processor needs to be translated into a physical address. It is simple process from mapping a virtual page to a physical page. Virtual pages of the first 16 segments bypass the TLB. The virtual address is identical to the physical address.

\begin{figure}[H]
\begin{center}
\begin{tikzpicture}[scale=0.9, transform shape]
        
  	\draw[help lines, gray!70, dashed] (0,0) grid(16,8);
  	
 	\def\vadrY{7}
 	\def\tlbY{4}
 	\def\padrY{0.5}
 	
 	% vAdr
  	\node[  tsRectangle, 
            minimum width=12cm,
            minimum height=1cm,
            fill=white] at (8,\vadrY) { };
                
  	\draw[-] (11,\vadrY - 0.5) -- (11, \vadrY + 0.5);
    \draw[-] (4,\vadrY - 0.5) -- (4, \vadrY + 0.5);
    
    \node at (13.75,\vadrY + 0.75) {0};
	\node at (11.25,\vadrY + 0.75) {13};
	\node at (10.75,\vadrY + 0.75) {12};
	\node at (4.25,\vadrY + 0.75) {51};
	\node at (3.75,\vadrY + 0.75) {52};
	\node at (2.25,\vadrY + 0.75) {63};           	
    \node at (3,\vadrY) {0};
    \node at (7,\vadrY) {Virtual page};
    \node at (12.5,\vadrY) {Page offset};
    
    %TLB
    \node[  tsRoundedRectangle, 
            minimum width=7cm,
            minimum height=2cm,
            fill=white] at (7.5,\tlbY) { TLB };

	% pAdr
  	\node[  tsRectangle, 
            minimum width=12cm,
            minimum height=1cm,
            fill=white] at (8,\padrY) { };
            
   	\draw[-] (11,\padrY - 0.5) -- (11, \padrY + 0.5);
    \draw[-] (6.5,\padrY - 0.5) -- (6.5, \padrY + 0.5);
    
    \node at (13.75,\padrY + 0.75) {0};
	\node at (11.25,\padrY + 0.75) {13};
	\node at (10.75,\padrY + 0.75) {14};
	\node at (6.75,\padrY + 0.75) {35};
	\node at (6.25,\padrY + 0.75) {36};
	\node at (2.25,\padrY + 0.75) {63};           	
    \node at (4.5,\padrY) {0};
    \node at (8.5,\padrY) {Physical page};
    \node at (12.5,\padrY) {Page offset};
    
    %arrows
    \draw[->, thick] (7.5, \vadrY - 0.5) -- (7.5, 5);
    \draw[<-, thick] (8.5, \padrY + 0.5) -- (8.5, 3);

  	                
\end{tikzpicture}
\end{center}
\caption{Address translation}
\label{fig:Address translation}
\end{figure}

\section*{Access Control}
\addcontentsline{toc}{section}{Access Control}

Twin-64 implements a protection model along two dimensions.  The first dimension is the \textbf{access control} and \textbf{privilege level} check. Each page is associated with a page type. There are read-only, read-write, execute and gateway pages. The page type is encoded in the AccType field. A value of zero is read-only access, a value of 1 is read-write access, a value of 2 is execute access and a value of 3 represents the gateway page access.

Each memory access is checked to be compatible with the allowed type of access set in the page descriptor. There are two privilege levels, user and supervisor mode. Access type and privilege level form the access control information field, which is checked for each instruction and data memory access. 

\begin{figure}[H]
\begin{center}
\begin{tikzpicture}[scale=0.9, transform shape]

	\def\leftofs{0}
        
  	\draw[help lines, gray!70, dashed] (0,0) grid(4,1.5);
  	
    \node[  tsRectangle, 
            minimum width=4cm,
            minimum height=1cm,  
            fill=white!50] (blocks) at (\leftofs + 2,1) { };
                
  	\draw[-] (\leftofs,0.5) -- (\leftofs,1.5);
    \draw[-] (\leftofs+2,0.5) -- (\leftofs+2,1.5);
    \draw[-] (\leftofs+3,0.5) -- (\leftofs+3,1.5);
  	
     	
    \def\flagOfs{0.125}
	\node at (\leftofs + 1, 1) {AcccType};
	\node at (\leftofs + 2.5, 1) {PL1};
	\node at (\leftofs + 3.5, 1) {PL2};
	
	\node at (\leftofs+0.5,0.25) {3};
	\node at (\leftofs+1.5,0.25) {2};
    \node at (\leftofs+2.5,0.25) {1};
   	\node at (\leftofs+3.5,0.25) {0};
   	 	
\end{tikzpicture}
\end{center}
\caption{Access control information field}
\label{fig:Access control information field}
\end{figure}

For read access the privilege level us be least as high as the PL1 field, for write access the privilege level must be as least as high as PL2. An exception are the execute and gateway pages, where a write access is only allowed for privilege code. For execute access the privilege level must be at least as high as PL1 and not higher than PL2. If the instruction privilege level is not within this bounds, a privilege violation trap is raised. If the page type does not match the instruction access type an access violation trap is raised. In both cases the instruction is aborted and a memory protection trap is raised.

\begin{longtable}{@{}p{0.2\linewidth}p{0.2\linewidth}p{0.2\linewidth}p{0.3\linewidth}@{}}
    \caption{Access type checking} \\
    \toprule
    \textbf{Type} & \textbf{read} & \textbf{write} & \textbf{execute} \\
    \midrule
    \endfirsthead
    \toprule
     \textbf{Type} & \textbf{read} & \textbf{write} & \textbf{execute} \\
    \midrule
    \endhead
    \midrule
    \multicolumn{2}{r}{\textit{Continued on next page}} \\
    \midrule
    \endfoot
    \bottomrule
    \endlastfoot
    \texttt{0 - read-only} & \texttt{PL <= PL1} & \texttt{not allowed} & \texttt{not allowed} \\
    \midrule 
    \texttt{1 - read-write} & \texttt{PL <= PL1} & \texttt{PL <= PL2} & \texttt{not allowed} \\
    \midrule 
    \texttt{2 - execute} & \texttt{PL <= PL1} & \texttt{PL == 0} & \texttt{PL2 <= PL <= PL1} \\
    \midrule 
    \texttt{3 - gateway} & \texttt{PL <= PL1} & \texttt{PL == 0} & \texttt{PL2 <= PL <= PL1} \\
    \midrule 
\end{longtable}


\section*{Access Protection}
\addcontentsline{toc}{section}{Access Protection}

The second dimension of protection is a \textbf{protection ID}. Twin-64 allows to record a set of up to eight protection Id values in the processor control registers \texttt{CR4} to \texttt{CR7} which are accessible to the currently executing process. The protection Id is identical to a segment Id. For each access to a segment that has the protection check bit set, one of the protection Id control registers must match the segment Id. If not, a protection violation trap is raised.

\begin{figure}[H]
\begin{center}
\begin{tikzpicture}[scale=0.9, transform shape]

	\def\leftofs{0}
        
  	\draw[help lines, gray!70, dashed] (0,0) grid(8,1.5);
  	
    \node[  tsRectangle, 
            minimum width=8cm,
            minimum height=1cm,  
            fill=white!50] (blocks) at (\leftofs + 4,1) { };
                
    \draw[-] (\leftofs+7.5,0.5) -- (\leftofs+7.5,1.5);
  	  	
   	\node at (\leftofs + 3.75, 1) {Protection Id};
	\node at (\leftofs + 7.75, 1) {W};
		
	\node at (\leftofs+0.5,0.25) {31};
    \node at (\leftofs+7.25,0.25) {1};
   	\node at (\leftofs+7.75,0.25) {0};
   	 	
\end{tikzpicture}
\end{center}
\caption{Protection Id}
\label{fig:Protection id}
\end{figure}


When protection Id checking is enabled, the eight protection identifiers are compared with the segment Id of the target address to validate an access. The rightmost bit of each of the eight protection Id registers is the write disable (W) bit. When the W bit is set, that protection Id cannot be used to grant write access. When the PSW P-bit is cleared, protection Id and write disable bit checking will not take place. 

Protection Id checking is typically used to form a grouping of access. A good example is the stack data segment, which is accessible at user privilege level. Since the CPU features a global virtual memory address space, every task could access a data stack of another task. With protection Id checking enabled and the segment Id loaded into the protection Id assigned control registers, only the corresponding user task is allowed to access the stack data segment with a matching protection Id. 

